
\chapter{Introdução}\label{cap:introducao}


O desenvolvimento de \textit{software} se torna mais complexo a cada dia e, com o surgimento constante de novas tecnologias e ferramentas, a quantidade de arquivos, métodos e estruturas envolvidos nos projetos cresce significativamente. Segundo \citeonline{Pressman2016}, essa complexidade crescente exige práticas cada vez mais organizadas para garantir qualidade, manutenibilidade e colaboração eficiente ao longo do ciclo de vida do \textit{software}.

Na indústria de desenvolvimento de jogos digitais, essa complexidade é ainda mais evidente, pois esse setor reúne programação e entretenimento. Criar um jogo é um trabalho que envolve diversos profissionais de diferentes áreas e, analogamente ao que afirma \citeonline{Schell2008}, pode-se destacar a atuação dos programadores, responsáveis por transformar código em mecânicas, e dos artistas sonoros e visuais, que dão vida a essas mecânicas. Por exemplo, o programador cria o personagem que corre, o artista visual a animação e o artista sonoro o som da corrida. Por consequência, o desenvolvimento de jogos caracteriza-se como um processo essencialmente colaborativo, no qual múltiplas equipes trabalham simultaneamente sobre os mesmos arquivos e sistemas.

Nesse sentido, quando se observa o cenário de colaboração, torna-se necessário o uso de ferramentas que permitam o trabalho conjunto de forma organizada e segura. Nesse contexto, destaca-se o Git, que, segundo \citeonline{Cunha2018}, é um sistema de controle de versões distribuído que permite a colaboração recíproca entre desenvolvedores, no qual cada integrante trabalha em seu ambiente local, ou repositório, e, após concluir uma tarefa, envia suas alterações para o repositório principal do projeto. 

Entretanto, apesar de sua ampla utilização e eficiência no controle de versões de arquivos textuais, como código-fonte, o Git apresenta limitações no gerenciamento de arquivos binários de grande porte, comuns no desenvolvimento de jogos, como imagens, modelos tridimensionais, áudios e vídeos. Para mitigar esse problema, utiliza-se o \gls{gitlfs}, uma extensão que armazena apenas o endereçamento desses arquivos no repositório principal, mantendo os dados reais em um servidor externo. Entretanto, essa abordagem não resolve plenamente os desafios relacionados ao versionamento eficiente, ao consumo de armazenamento e ao controle detalhado das alterações em arquivos binários.

Diante desse contexto, observa-se que os mecanismos tradicionais de controle de versões apresentam limitações significativas no gerenciamento eficiente de arquivos binários de grande porte. Embora soluções como o \gls{gitlfs} reduzam parcialmente o impacto desses arquivos nos repositórios, tais abordagens ainda tratam os dados binários como unidades indivisíveis, resultando no armazenamento completo de novas versões mesmo quando apenas pequenas partes do arquivo foram modificadas. Nesse cenário, surge a necessidade de investigar abordagens alternativas capazes de lidar de forma mais eficiente com o versionamento desses dados e analisar também quais os melhores dados para tal.

Desta forma, este trabalho busca responder à seguinte questão de pesquisa: Como técnicas de armazenamento diferencial por blocos, baseadas nos princípios de \textit{Copy-on-Write} e \textit{Redirect-on-Write}, podem ser aplicadas ao versionamento de arquivos binários para reduzir custos de armazenamento e melhorar o desempenho em ambientes colaborativos de desenvolvimento de software?

\section{Objetivos}\label{sec:objetivos}

Os objetivos do trabalho se dividem em objetivo geral e objetivos específicos.

\subsection{Objetivo geral}\label{subsec:objetivoGeral}

Desenvolver um novo mecanismo de versionamento de arquivos binários fundamentado em arquiteturas baseadas nos princípios de \gls{cow} e \gls{row}, visando otimizar o armazenamento, o desempenho e o controle das alterações, especialmente no contexto do desenvolvimento de jogos digitais em três dimensões.

\subsection{Objetivos específicos}\label{subsec:objetivosEspecificos}

\begin{itemize}
    \item Analisar as limitações dos mecanismos tradicionais de versionamento no tratamento de arquivos binários, incluindo a identificação de cenários nos quais a divisão em blocos seja viável como objeto de estudo.
    
    \item Investigar a aplicação das técnicas \gls{cow} e \gls{row} no contexto de armazenamento diferencial por blocos, considerando sistemas de arquivos como o \gls{btrfs}  e o \gls{zfs}, bem como seu funcionamento no Kernel do Linux e sua aplicabilidade ao problema proposto.
    
    \item Propor uma abordagem para o versionamento de arquivos binários baseada em armazenamento diferencial por blocos, avaliando sua eficiência em termos de consumo de armazenamento e desempenho, por meio de análise comparativa com abordagens utilizadas no mercado.
    
\end{itemize}


%Os objetivos específicos são opcionais, ou seja, somente devem ser apresentados se caracterizarem resultados parciais gerados a partir do objetivo geral, os quais sejam considerados úteis para a comunidade acadêmica, para a sociedade ou para o ambiente profissional. Uma observação importante é que os resultados sejam passíveis de comprovação, ou seja, se o objetivo for: “Oferecer agilidade e confiabilidade aos processos gerenciais da empresa”, significa que o trabalho deverá realizar testes com relação a esses atributos, cujos resultados deverão ser apresentados nas discussões do trabalho.

%Destaca-se que os objetivos específicos não incluem as etapas do processo de desenvolvimento de software (realizar a modelagem, a análise, o projeto...) ou outras atividades necessárias para alcançar o objetivo geral, como, estudar as tecnologias necessárias para modelagem e implementação do sistema. Dentre as exceções estão a realização de estudos, procedimentos, métodos e técnicas considerados inéditos e de relevância para outros trabalhos a serem realizados na mesma área. Contudo, o resultado deste estudo deve ser documentado de forma que seja conhecimento disponibilizado para quem lê o trabalho.

\section{Justificativa}\label{sec:justificativa}

Ao tratar de arquivos binários no contexto de ambientes de produção, é possível citar exemplos como arquivos de áudio, imagens em alta resolução, modelos tridimensionais e até mesmo modelos de inteligência artificial. Esses arquivos, quando armazenados em repositórios que seguem fluxos tradicionais de versionamento, os quais mantêm o histórico completo das alterações, apresentam limitações significativas. Tal cenário ocorre devido à impossibilidade de comparar esses dados de forma eficiente por meio de operações de diferença, \textit{diff}, além da dificuldade de compactá-los adequadamente, o que pode comprometer o desempenho e o armazenamento do repositório ao longo do tempo.

\gls{gitlfs} surge como uma solução paliativa para esse problema, pois substitui arquivos binários de grande porte por arquivos de ponteiro que referenciam servidores externos, nos quais os dados reais são armazenados \cite{Chen2025}. Apesar de essa abordagem reduzir o tamanho do repositório principal, ela não resolve de maneira estrutural os desafios relacionados ao versionamento eficiente desses arquivos.

O principal problema reside no fato de que os arquivos binários passam a ser tratados como “caixas pretas”, cujo conteúdo interno não pode ser analisado e comparado. Por consequência, qualquer modificação, mesmo que mínima, resulta no envio e armazenamento de uma nova versão completa do arquivo, ocasionando desperdício significativo de espaço em disco, aumento no tráfego de dados e impacto negativo no desempenho dos sistemas de versionamento.

Diante desse cenário, torna-se necessário investigar alternativas capazes de evitar a substituição integral desses binários, permitindo o registro apenas das porções efetivamente modificadas dos arquivos. Nesse contexto, propõe-se a utilização de sistemas de versionamento e armazenamento fundamentados nos princípios de \gls{cow} e \gls{row}, amplamente utilizadas em sistemas de armazenamento para a criação de \textit{snapshots} diferenciais, conforme apresentado por \cite{Xiao2009}. Essas abordagens possibilitam a preservação de versões anteriores por meio do armazenamento seletivo de blocos alterados, reduzindo significativamente o consumo de espaço e os custos computacionais associados às operações de escrita.

Nesse sentido, a aplicação desses conceitos ao contexto do versionamento de arquivos binários apresenta-se como uma alternativa promissora para superar as limitações dos mecanismos tradicionais, possibilitando maior eficiência no armazenamento, melhoria no desempenho e controle mais refinado das alterações, especialmente em ambientes colaborativos de grande escala, como os encontrados na indústria de desenvolvimento de jogos digitais, onde os arquivos binários podem mudar a mesma medida que os arquivos de texto.
